% Part: first-order-logic
% Chapter: completeness
% Section: completeness-theorem

\documentclass[../../../include/open-logic-section]{subfiles}

\begin{document}

\iftag{FOL}
      {\olfileid{fol}{com}{cth}}
      {\olfileid{pl}{com}{cth}}

\olsection{مبرهنة الاكتمال}

\begin{explain}
لنجمع نتائجنا، فنصل إلى مبرهنة الاكتمال.
\end{explain}

\begin{thm}[مبرهنة الاكتمال]
\ollabel{thm:completeness}
لتكن $\Gamma$ مجموعة من ال!!{sentence}s. إذا كانت $\Gamma$ متسقة، فهي
قابلة للإرضاء.
\end{thm}

\begin{proof}
افترض أن $\Gamma$ متسقة.\iftag{FOL}{ بحسب \olref[hen]{lem:henkin}،
  توجد مجموعة متسقة مشبعة $\Gamma' \supseteq \Gamma$.}{} وبحسب
\olref[lin]{lem:lindenbaum} توجد
$\Gamma^* \supseteq \iftag{FOL}{\Gamma'}{\Gamma}$ متسقة و!!{complete}.
\iftag{FOL}{
  لما كان $\Gamma' \subseteq \Gamma^*$، فإن $\Gamma^*$ تحتوي، لكل
  !!{formula}~$!A(x)$، !!a{sentence} من الصورة
  \iftag{prvEx}
        {$\lexists[x][!A(x)] \lif !A(c)$}
        {$\lnot\lforall[x][!A(x)] \lif \lnot !A(c)$}
  ومن ثم تكون $\Gamma^*$ مشبعة. وإذا كانت $\Gamma$ لا تحتوي~$\eq$،
  فبحسب}{بحسب}
\olref[mod]{lem:truth},
$\iftag{FOL}{\Sat{M(\Gamma^*)}}{\pSat{v(\Gamma^*)}}{!A}$ إذا وفقط إذا
كان $!A \in \Gamma^*$. وينتج من ذلك بوجه خاص أنه لكل $!A \in
\Gamma$، يتحقق $\iftag{FOL}{\Sat{M(\Gamma^*)}}{\pSat{v(\Gamma^*)}}{!A}$،
فتكون $\Gamma$ قابلة للإرضاء.\iftag{FOL}{
  وإذا كانت $\Gamma$ تحتوي~$\eq$، فإن \olref[ide]{lem:truth} تعطينا،
  لكل ال!!{sentence}s~$!A$، أن $\Sat{\equivclass{M}{\approx}}{!A}$ إذا
  وفقط إذا كان $!A \in \Gamma^*$. وبوجه خاص، يتحقق
  $\Sat{\equivclass{M}{\approx}}{!A}$ لكل $!A \in \Gamma$، فتكون
  $\Gamma$ قابلة للإرضاء.}{}
\end{proof}

\begin{cor}[مبرهنة الاكتمال، الصيغة الثانية]
\ollabel{cor:completeness}
لكل $\Gamma$ ولكل !!{sentence}~$!A$: إذا كان $\Gamma \Entails !A$،
فإن $\Gamma \Proves !A$.
\end{cor}

\begin{proof}
لاحظ أن كلًا من $\Gamma$ في \olref{cor:completeness} و
\olref{thm:completeness} مكمّم كليًا. ولكي لا نخلط بينهما، لنعِد صياغة
\olref{thm:completeness} باستعمال متغير مختلف: لكل مجموعة من
ال!!{sentence}s~$\Delta$، إذا كانت $\Delta$ متسقة فهي قابلة للإرضاء.
وبالمقابل المعاكس، إذا لم تكن $\Delta$ قابلة للإرضاء فهي غير متسقة.
وسنستعمل ذلك لبرهان النتيجة.

افترض أن $\Gamma \Entails !A$. عندئذ تكون $\Gamma \cup \{\lnot !A\}$
غير قابلة للإرضاء بحسب \olref[syn][sem]{prop:entails-unsat}. وإذا اتخذنا
$\Gamma \cup \{\lnot !A\}$ هي $\Delta$، أعطتنا الصيغة السابقة من
\olref{thm:completeness} أن $\Gamma \cup \{\lnot !A\}$ غير متسقة.
وبحسب
\iftag{FOL}{%
  \tagrefs{prfAX/{fol:axd:prv:prop:prov-incons},
    prfSC/{fol:seq:prv:prop:prov-incons},
    prfND/{fol:ntd:prv:prop:prov-incons},
    prfTab/{fol:tab:prv:prop:prov-incons}}}{%
  \tagrefs{prfAX/{pl:axd:prv:prop:prov-incons},
    prfSC/{pl:seq:prv:prop:prov-incons},
    prfND/{pl:ntd:prv:prop:prov-incons},
    prfTab/{pl:tab:prv:prop:prov-incons}}},
نحصل على $\Gamma \Proves !A$.
\end{proof}

\tagprob{FOL}
\begin{prob}
استعمل \olref[fol][com][cth]{cor:completeness} لبرهان
\olref[fol][com][cth]{thm:completeness}، مبيّنًا بذلك أن صيغتي مبرهنة
الاكتمال متكافئتان.
\end{prob}
\tagendprob

\tagprob{notFOL}
\begin{prob}
استعمل \olref[pl][com][cth]{cor:completeness} لبرهان
\olref[pl][com][cth]{thm:completeness}، مبيّنًا بذلك أن صيغتي مبرهنة
الاكتمال متكافئتان.
\end{prob}
\tagendprob

\tagprob{FOL}
\begin{prob}
لكي يكون نسق !!a{derivation} كاملًا، يجب أن تكون قواعده قوية بما يكفي
لإثبات أن كل مجموعة غير قابلة للإرضاء غير متسقة. فأي قواعد
ال!!{derivation} كانت لازمة لبرهان الاكتمال؟ وهل توجد من هذه القواعد
قاعدة لم تُستعمل في أي موضع من البرهان؟ وللإجابة عن هذين السؤالين، ضع
قائمة أو مخططًا يبيّن أي قواعد ال!!{derivation} استُعملت في أي النتائج
المؤدية إلى برهان \olref[fol][com][cth]{thm:completeness}. واحرص على
تسجيل كل استعمال ضمني للقواعد في هذه البراهين.
\end{prob}
\tagendprob

\tagprob{notFOL}
\begin{prob}
لكي يكون نسق !!a{derivation} كاملًا، يجب أن تكون قواعده قوية بما يكفي
لإثبات أن كل مجموعة غير قابلة للإرضاء غير متسقة. فأي قواعد
ال!!{derivation} كانت لازمة لبرهان الاكتمال؟ وهل توجد من هذه القواعد
قاعدة لم تُستعمل في أي موضع من البرهان؟ وللإجابة عن هذين السؤالين، ضع
قائمة أو مخططًا يبيّن أي قواعد ال!!{derivation} استُعملت في أي النتائج
المؤدية إلى برهان \olref[pl][com][cth]{thm:completeness}. واحرص على
تسجيل كل استعمال ضمني للقواعد في هذه البراهين.
\end{prob}
\tagendprob

\end{document}
